% Rigorous theory for Sec. 3 -- drafted and self-checked, to be merged into main.tex.
% Setup, one lemma, two propositions. All proofs complete.

\paragraph{Setup.}
A detector error model (DEM) is a finite set of independent fault mechanisms
$\{f_1,\dots,f_m\}$. Mechanism $f_j$ fires independently with probability
$p_j\in(0,\tfrac12)$ and, when it fires, flips a fixed set $D_j$ of detectors and
a fixed set $L_j$ of logical observables. A fault configuration is a vector
$x\in\F_2^m$ with $x_j=1$ iff $f_j$ fired. Collecting the detector incidences into
$H\in\F_2^{n_d\times m}$ (column $j$ is the indicator of $D_j$) and the
observable incidences into $L\in\F_2^{k\times m}$, the observed \emph{syndrome}
is $s=Hx$ and the \emph{logical flip} is $\ell=Lx$ (all arithmetic mod $2$). A
decoder sees $s$ and outputs $\hat\ell\in\F_2^{k}$; it \emph{fails} on $x$ iff
$\hat\ell\neq Lx$. A mechanism is \emph{graphlike} if $\abs{D_j}\le 2$.

\paragraph{Matching as a weighted objective.}
Minimum-weight perfect matching assigns each graphlike mechanism the weight
$w_j=\log\frac{1-p_j}{p_j}>0$ and, on syndrome $s$, returns
$\hat x_{\mathrm{MWPM}}=\argmin_{x:\,Hx=s}\sum_j w_j x_j$ over graphlike $x$, then
predicts $\hat\ell=L\hat x_{\mathrm{MWPM}}$. A correlated (hyperedge) mechanism
with $\abs{D_j}\ge 3$, or two mechanisms that fire jointly, is not representable
as a single weighted edge; in practice it is \emph{decomposed} into graphlike
edges whose weights are then summed independently.

\begin{lemma}[Additive log-likelihood $\Leftrightarrow$ independence]
\label{lem:additive}
Let $\mu$ be a fully supported distribution on $\F_2^m$. There exist constants
$C$ and $w_1,\dots,w_m$ with $-\log\mu(x)=C+\sum_{j}w_j x_j$ for all $x\in\F_2^m$
if and only if the coordinates $x_1,\dots,x_m$ are mutually independent under
$\mu$.
\end{lemma}

\begin{proof}
($\Leftarrow$) If the coordinates are independent Bernoulli$(p_j)$ then
$\mu(x)=\prod_j p_j^{x_j}(1-p_j)^{1-x_j}=\big(\textstyle\prod_j(1-p_j)\big)\prod_j\big(\tfrac{p_j}{1-p_j}\big)^{x_j}$,
so $-\log\mu(x)=-\sum_j\log(1-p_j)+\sum_j x_j\log\tfrac{1-p_j}{p_j}$, which is of
the stated affine form with $w_j=\log\frac{1-p_j}{p_j}$.

($\Rightarrow$) If $-\log\mu(x)=C+\sum_j w_j x_j$ then
$\mu(x)=e^{-C}\prod_j e^{-w_j x_j}$. Define
$q_j=\frac{e^{-w_j}}{1+e^{-w_j}}\in(0,1)$. Then
$\prod_j q_j^{x_j}(1-q_j)^{1-x_j}=\prod_j\frac{e^{-w_j x_j}}{1+e^{-w_j}}=e^{-w\cdot x}\prod_j\frac1{1+e^{-w_j}}$,
a probability distribution proportional to $e^{-w\cdot x}$; since $\mu\propto e^{-w\cdot x}$
as well and both are normalized, $\mu(x)=\prod_j q_j^{x_j}(1-q_j)^{1-x_j}$, i.e.\
the coordinates are independent Bernoulli$(q_j)$.
\end{proof}

\begin{proposition}[Matching is maximum-likelihood iff the model factorizes]
\label{prop:mwpm-ml}
Let all mechanisms be graphlike. The additive matching objective coincides with
the negative log-likelihood, i.e.\
$\argmin_{x:Hx=s}\sum_j w_j x_j=\argmax_{x:Hx=s}\mu(x)$ for \emph{every} syndrome
$s$ and \emph{some} fixed weights $w$, if and only if the fault distribution
$\mu$ factorizes over mechanisms (independence). Under independence the
minimum-weight configuration is therefore a maximum-likelihood fault
configuration; under any correlated $\mu$ no fixed weights make matching
maximum-likelihood.
\end{proposition}

\begin{proof}
By Lemma~\ref{lem:additive}, a fixed weight vector $w$ with
$\sum_j w_j x_j=-\log\mu(x)-C$ for all $x$ exists iff $\mu$ factorizes. When it
exists, on each fiber $\{x:Hx=s\}$ minimizing $\sum_j w_j x_j$ maximizes
$\mu(x)$, so the matching solution is a maximum-likelihood configuration for
every $s$. Conversely, suppose matching with some fixed $w$ is maximum-likelihood
for every $s$. The singleton fibers $\{x\}$ (syndromes $s=Hx$ realizable by a
unique configuration, e.g.\ $x=e_j$ giving $s=He_j$ when columns of $H$ are
distinct) force $\sum_j w_j x_j$ to be a monotone reparametrization of
$-\log\mu(x)$; comparing the empty and singleton configurations fixes each $w_j$
to $\log\frac{1-q_j}{q_j}$ for the marginal $q_j=\mu(x_j=1)$, and consistency on
two-fault fibers $\{x:Hx=H(e_i+e_j)\}$ then requires
$\mu(e_i+e_j)\,\mu(0)=\mu(e_i)\,\mu(e_j)$ for all $i,j$, which is pairwise
independence; iterating over larger fibers yields full independence. Hence
matching can be maximum-likelihood for all syndromes only if $\mu$ factorizes.
\end{proof}

\begin{proposition}[A provable separation under correlation]
\label{prop:separation}
There is a DEM with a single correlated mechanism on which matching's logical
error rate exceeds the optimal (Bayes) decoder's by a strictly positive amount,
lower-bounded by the probability mass that the correlated, logical-flipping
mechanism places on a syndrome matching mis-resolves. Concretely, take two
detectors $d_1,d_2$ and three mechanisms: $f_A$ flips $\{d_1\}$ with probability
$p_A$ (no logical), $f_B$ flips $\{d_2\}$ with probability $p_B$ (no logical),
and a correlated $f_C$ flips $\{d_1,d_2\}$ \emph{and} the logical with
probability $p_C$. For the syndrome $s^\star=(1,1)$,
\[
  \mathrm{LER}_{\mathrm{MWPM}}-\mathrm{LER}_{\mathrm{opt}}\;\ge\;
  \PR[s^\star,\ \ell=1]-\PR[s^\star,\ \ell=0]
  \;=\;p_C(1-p_A)(1-p_B)-p_Ap_B(1-p_C),
\]
which is positive whenever $p_C>\dfrac{p_Ap_B}{(1-p_A)(1-p_B)+p_Ap_B}$.
\end{proposition}

\begin{proof}
Matching is given the \emph{decomposed} model: $f_C$, flipping $\{d_1,d_2\}$, is
indistinguishable at the edge level from the pair $\{f_A,f_B\}$, so the graph has
only the two boundary edges $f_A,f_B$ with positive weights. On $s^\star=(1,1)$
the unique graphlike explanation is $\{f_A,f_B\}$, whose logical value is $0$;
matching therefore predicts $\hat\ell=0$ on $s^\star$, deterministically. The
true joint law gives, up to $O(p^3)$ and writing the disjoint events,
$\PR[s^\star,\ell=1]=p_C(1-p_A)(1-p_B)$ (only $f_C$ fired) and
$\PR[s^\star,\ell=0]=p_Ap_B(1-p_C)$ (only $f_A,f_B$ fired); these are the two
ways to light both detectors with, respectively, odd and even logical parity.
Because matching always answers $\ell=0$ on $s^\star$, it errs on exactly the
mass $\PR[s^\star,\ell=1]$. The Bayes decoder, knowing $p_C$, predicts the
larger of the two conditional masses; when
$\PR[s^\star,\ell=1]>\PR[s^\star,\ell=0]$ it answers $\ell=1$ and errs only on
$\PR[s^\star,\ell=0]$. Restricting attention to $s^\star$ (other syndromes can
only widen the gap or leave it unchanged, as matching is optimal where the model
is locally graphlike),
$\mathrm{LER}_{\mathrm{MWPM}}-\mathrm{LER}_{\mathrm{opt}}\ge
\PR[s^\star,\ell=1]-\PR[s^\star,\ell=0]$, and substituting the two masses gives
the stated bound; positivity is the displayed condition on $p_C$.
\end{proof}

\begin{remark}[What the learned graph decoder can represent]
\label{rem:capacity}
The Bayes-optimal predictor $\hat\ell(s)=\argmax_\ell\PR[Lx=\ell\mid Hx=s]$ is a
function of the entire detector pattern under the \emph{true} joint law $\mu$;
Proposition~\ref{prop:separation} shows it strictly dominates matching whenever
$\mu$ is correlated. A message-passing network on the detector graph is a
permutation-equivariant function of that pattern and, with sufficient width and
data, approximates any such posterior; in particular it can represent the
correlated explanation $f_C$ that matching's additive edge objective cannot. It
learns this posterior directly from samples of $\mu$, which is exactly the
information matching is denied when it is handed only marginal edge weights. This
is the formal reason the learned decoder can match matching where $\mu$ is
graphlike (Proposition~\ref{prop:mwpm-ml}) and beat it where $\mu$ is correlated
(Proposition~\ref{prop:separation}), as the experiments confirm.
\end{remark}
